\documentclass[a4paper]{article}


%-----------------------PACKAGES----------------------------------------


\usepackage[brazilian]{babel}
%\usepackage[latin1]{inputenc}
\usepackage[utf8]{inputenc}
\usepackage{lipsum}
\usepackage{amsmath}
\usepackage{mathtools}
\usepackage{amsthm}
\usepackage{amsfonts}
\usepackage{multirow}
\usepackage{wrapfig}
%\usepackage{natbib}
\usepackage{footnote}
\usepackage{booktabs, ctable}
\usepackage{enumerate}
\usepackage{array}
\usepackage{placeins}
\usepackage{float}
\usepackage{caption}
\usepackage{subcaption}
\usepackage{amssymb}
\usepackage{indentfirst}
\usepackage{cleveref}
\usepackage{template-unb}
\graphicspath{{img/}}


%-----------------------MACROS-------------------------------------------
\newcommand{\vars}[1]{\mathcal{V}ars(#1)}
\newcommand{\pos}[1]{\mathcal{P}os(#1)}
\newcommand{\terms}[2]{\mathcal{T}(#1, #2)}
\newcommand{\termpos}[1]{|_{#1}}
\newcommand{\subpos}[2]{[#1]_{#2}}
\newcommand{\dom}[1]{\mathcal{D}om(#1)}
\newcommand{\parared}[2]{\rightrightarrows_{#1}^{#2}}





%-----------------------THEOREMS----------------------------------------
\theoremstyle{definition}
\newtheorem{definition}{Definição}
\theoremstyle{plain}
\newtheorem{proposition}[definition]{Proposição}
\newtheorem{theorem}[definition]{Teorema}
\newtheorem{lemma}[definition]{Lema}
\newtheorem{corollary}[definition]{Corolário}
\newtheorem{fact}[definition]{Fato}
\crefname{definition}{definição}{definições}
\Crefname{definition}{Definição}{Definições}
\crefname{proposition}{proposição}{proposições}
\Crefname{proposition}{Proposição}{Proposições}
\crefname{theorem}{teorema}{teoremas}
\Crefname{theorem}{Teorema}{Teoremas}
\crefname{lemma}{lema}{lemas}
\Crefname{lemma}{Lema}{Lemas}
\crefname{corollary}{corolário}{corolários}
\Crefname{corollary}{Corolário}{Corolários}
\crefname{fact}{fato}{fatos}
\Crefname{fact}{Fato}{Fatos}
\renewenvironment{proof}{{\bfseries \noindent Demonstração:}}{\qed}
\renewcommand{\qedsymbol}{$\blacksquare$}

%-----------------------DOCUMENT----------------------------------------



\begin{document}
	\Titulo{Avaliação 1}{Formalismos Matemáticos da Computação}{Guilherme Borges Brandão}

\vspace{-0.5cm}
\begin{abstract}
		Este documento tem como objetivo providenciar as provas formais do Teorema 6.3.10 e dos Lemas 6.4.4 e 6.4.8 do livro {\it Term Rewriting and All That}, de Franz Baader e Tobias Nipkow, como parte da Avaliação 1 da disciplina {\it Formalismos Matemáticos da Computação}.
	\end{abstract}

	\begin{definition}
		Sejam \(\Sigma\) uma assinatura, \(X\) um conjunto de variáveis disjuntas de \(\Sigma\) e \(s,t\in \terms{\Sigma}{X}\).
		\begin{enumerate}
			\item O {\em conjunto de posições} de \(s\) é definido recursivamente como
				\[\pos{s} := \begin{cases}
					\{\epsilon\} & \text{se } s \in X \\
					\{\epsilon\} \cup \bigcup_{i=1}^{n}\{i p \mid p \in \pos{s_i}\} & \text{se } s = f(s_1, \ldots, s_n)
				\end{cases}\]
				onde \(\epsilon\) é a {\em posição raiz} de \(s\), e a função ou símbolo de variável nessa posição é chamado de símbolo raíz. A {\em ordem de prefixo} definida como 
					\[p \leq q \text{ sse exuste } p' \text{ such tal que } pp' = q\]
				é uma ordem parcial sobre \(\pos{s}\). A posição \(p\) é dita {\em acima} de \(q\) se \(p < q\) e {\em abaixo} de \(q\) se \(p > q\). Dizemos que duas posições $p$ e $q$ de um termo $t$ são paralelas se não existe uma posição $r$ tal que \(r < p\) e \(r < q\) e denotamos por \(p \parallel q\).
			\item O {\em tamanho} de \(s\) é definido como \(|s| := |\pos{s}|\).
			\item O {\em subtermo} de \(s\) na posição \(p\) é definido recursivamente como
				\[s\termpos{p} := \begin{cases}
					s & \text{se } p = \epsilon \\
					s_i\termpos{p'} & \text{se } s = f(s_1, \ldots, s_n) \text{ e } p = ip' \text{ para algum } i \in \{1, \ldots, n\}.
				\end{cases}\]
			\item Para \(p\in \pos{s}\), denotamos por \(s\subpos{t}{p}\) o termo obtido substituindo o subtermo de \(s\) na posição \(p\) pelo termo \(t\). Formalmente, temos
				\[s\subpos{t}{p} := \begin{cases}
					t & \text{se } p = \epsilon \\
					f(s_1, \ldots, s_{i-1}, s_i\subpos{t}{p'}, s_{i+1}, \ldots, s_n) & \text{se } s = f(s_1, \ldots, s_n).
				\end{cases}\]
			\item Por \(\vars{s}\), denotamos o conjunto de variáveis que ocorrem em \(s\).
				\[\vars{s} := \{x \in X \mid \exists p \in \pos{s} \text{ tal que } s\termpos{p} = x\}\]
			Dizemos que \(p\in\pos{s}\) é uma {\em posição de variável} se \(s\termpos{p}\) é uma variável.
		\end{enumerate}
	\end{definition}

	

	\begin{lemma}
		\label{lemma:paralelo-subst1}
		Sejam \(s, t, r\) termos e \(p,q\in \pos{s}\) palavras sobre os inteiros positivos. 
		\begin{enumerate}
			\item Se \(pq\in \pos{s}\), então \(s\pos{pq} = (s\pos{p})\pos{q}\).
			\item Se \(p\in\pos{s}\) e \(q\in\pos{t}\), então
				\[\begin{aligned}
					(s\subpos{t}{p})\termpos{pq} & = t \termpos{q} \\
				    (s\subpos{t}{p})\subpos{r}{pq} & = s\subpos{t\subpos{r}{q}}{p}
				\end{aligned}\]
			\item Se \(pq\in \pos{s}\), então
				\[\begin{aligned}
					(s\subpos{t}{pq})\termpos{p} & = (s\termpos{p})\subpos{t}{q}\\
				    (s\subpos{t}{pq})\subpos{r}{p} & = s\subpos{r}{p}
				\end{aligned}\]
			\item Se \(p \parallel q\), então
				\[\begin{aligned}
					(s\subpos{t}{p})\termpos{q} & = s\termpos{q} \\
				    (s\subpos{t}{p})\subpos{r}{q} & = (s\subpos{r}{q})\subpos{t}{p}
				\end{aligned}\]
		\end{enumerate}
	\end{lemma}
	\begin{proof}
	A prova é por indução na posição \(p\) seguida de uma análise de casos na estrutura de \(s\).
	\end{proof}

	\begin{definition}
		Sejam \(\Sigma\) uma assinatura e \(V\) um conjunto enumerável de varíaveis. Uma {\em \(\terms{\Sigma}{V}\)-substituição} é uma função \(\sigma \colon V \to \terms{\Sigma}{V}\) tal que \(\{x \in V \mid \sigma(x) \neq x\}\) é finito. Tal conjunto é chamado de {\em domínio} de \(\sigma\) e denotado por \(\dom{\sigma}\). Se \(\dom{\sigma}=\{x_1, \ldots, x_n\}\), escrevemos \(\sigma = \{x_1 \mapsto \sigma(x_1), \ldots, x_n \mapsto \sigma(x_n)\}\).
	\end{definition}

	\begin{definition}
		Um sistema de reescrita de termos é dito ortogonal se ele é linear à esquerda e não possui pares críticos.
	\end{definition}

	\begin{definition}
		Dizemos que um conjunto de posições \(P\) de um termo \(t\) é paralelo se \(p \parallel q\) ou \(p = q\) para todo \(p, q \in P\). Seja \(P\) um conjunto de posições paralelo. Dado \(t_p\) para cada \(p \in P\) definimos a notação
			\[s\subpos{t_p}{p\in P} := s\subpos{t_{p_1}}{p_1}\ldots \subpos{t_{p_n}}{p_n}.\]
		Como cada \(p_i\) é paralelo entre si, a ordem das substituições não importa.

		Se para cada \(p \in P\) temos \(l_p \rightarrow r_p\in R\) e uma substituição \(\sigma_p\) onde \(s\termpos{p} = \sigma_p l_p\), escrevemos
			\[s \parared{R}{P} s\subpos{\sigma_p r_p}{p\in P}\]
		e chamamos isso de um passo de redução paralela. As anotações de \(R\) e \(P\) podem ser omitidas quando não houver ambiguidade.
	\end{definition}

	Uma consequência imediata da definição acima é

	\begin{fact}
	\(\rightarrow \ \subset \ \parared{}{} \ \subset \ \rightarrow^*\)
	\end{fact}
	\begin{proof}
		\begin{enumerate}
			\item Primeiro mostremos que \(\rightarrow \ \subset \ \parared{}{}\). Suponha \(s\rightarrow t\), o objetivo é mostrar que \(s \parared{}{} t\). Como \(s\rightarrow t\), existe uma posição \(p\) tal que \(s\termpos{p} = \sigma l\) e \(t = s\subpos{\sigma r}{p}\) para alguma regra \(l \rightarrow r \in R\) e substituição \(\sigma\). Temos que conjunto \(\{p\}\) é paralelo. Portanto, temos
				\[s \parared{R}{\{p\}} s\subpos{\sigma r}{p} = t.\]

			\item Agora, vamos mostrar que \(\parared{}{} \ \subset \ \rightarrow^*\). Suponha \(s \parared{R}{P} t\). Disso segue que existe um conjunto de posições paralelo \(P = \{p_1, p_2, \ldots, p_n\}\) e para cada \(p_i \in P\) existe uma regra \(l_i \rightarrow r_i \in R\) e uma substituição \(\sigma_i\) tal que \(s\termpos{p_i} = \sigma_i l_i\), com 
				\[t = s\subpos{\sigma_{1} r_{1}}{p_1}\ldots \subpos{\sigma_{n} r_{n}}{p_n}.\]
			Como \(P\) é paralelo, cada redução unitária pode ser realizada em qualquer ordem, sem alterar os subtermos das outras posições (\Cref{lemma:paralelo-subst1}). Logo, temos a seguinte cadeia de reduções:
				\begin{align*}
					s &= s\subpos{\sigma_{1} l_{1}}{p_1}\ldots \subpos{\sigma_{n} l_{n}}{p_n} \\
					&\rightarrow s\subpos{\sigma_{1} r_{1}}{p_1}\subpos{\sigma_{2} l_{2}}{p_2}\ldots \subpos{\sigma_{n} l_{n}}{p_n} \\
					&\rightarrow s\subpos{\sigma_{1} r_{1}}{p_1}\subpos{\sigma_{2} r_{2}}{p_2}\ldots \subpos{\sigma_{n} l_{n}}{p_n} \\
					&\vdots \\
					&\rightarrow s\subpos{\sigma_{1} r_{1}}{p_1}\ldots \subpos{\sigma_{n} r_{n}}{p_n} = t.
				\end{align*}
			Ou seja, \(s \rightarrow^* t\).
		\end{enumerate}
	\end{proof}



\end{document}